\section{The landscape}

\begin{frame}{Two design loops, one pattern}
  \begin{columns}[T,onlytextwidth]
    \column{0.48\textwidth}
      \textbf{Analog loop}\\[0.3em]
      \begin{itemize}\small
        \item \term{Inputs}: spec (GBW, PM, noise), topology choice.
        \item \term{Evaluator}: ngspice + gm/ID LUTs.
        \item \term{Knobs}: W/L, bias current, mirror ratios.
        \item \term{FoM}: meets spec? area? power?
      \end{itemize}
    \column{0.48\textwidth}
      \textbf{Digital loop}\\[0.3em]
      \begin{itemize}\small
        \item \term{Inputs}: RTL, clock, die budget.
        \item \term{Evaluator}: LibreLane (Yosys + OpenROAD + KLayout).
        \item \term{Knobs}: density, clock period, PDN pitch.
        \item \term{FoM}: timing slack, area, power, DRC/LVS green.
      \end{itemize}
  \end{columns}
  \vspace{0.8em}
  \begin{center}
    \textit{Both loops share the pattern: propose $\rightarrow$ simulate
    $\rightarrow$ score $\rightarrow$ repeat. The \cmd{eda-agents} repo
    implements it twice, with the same autoresearch core.}
  \end{center}
  \note{%
    Core thesis. Analog and digital feel like different worlds -- one is
    Verilog and flows, the other is W/L ratios and SPICE. But at the
    agent level they are structurally identical: a proposal engine (the
    LLM) emits candidate parameter sets, an evaluator measures them, a
    FoM turns those measurements into a scalar, and a loop repeats.\par
    What differs is the per-evaluation cost. Analog: seconds (ngspice on
    a few transistors). Digital: minutes (LibreLane through routing +
    signoff). The budgets we spend and the search strategies we pick
    follow from that asymmetry.}
\end{frame}

\begin{frame}{What \cmd{eda-agents} ships today}
  \begin{center}
  \begin{tikzpicture}[node distance=0.5cm]
    \node[agent] (atr)  {analog-topology\\recommender};
    \node[agent, right=of atr] (asa) {analog-sizing\\advisor};
    \node[agent3, right=of asa] (gfa) {gf180-docker\\analog};
    \node[agent2, below=0.8cm of atr] (dtb) {digital-testbench\\author};
    \node[agent2, right=of dtb] (gfd) {gf180-docker\\digital};
    \node[agent, right=of gfd] (llx) {librelane\\expert};

    \node[draw=muted, dashed, rounded corners=2pt, font=\scriptsize,
          minimum height=0.55cm, minimum width=11.5cm]
         (mcp) at ($(asa.north)!0.5!(gfd.south) + (1.6cm, 1.2cm)$) {};
  \end{tikzpicture}
  \end{center}
  \vspace{0.2em}
  \footnotesize
  \textbf{Analog agents} (blue): recommend a topology, size it with gm/ID,
  harden on GF180 inside Docker.\\
  \textbf{Digital agents} (green): author cocotb TB, drive RTL-to-GDS on
  GF180 inside Docker.\\
  \textbf{Shared infra} (orange/muted): LibreLane expert + MCP server.
  \note{%
    This diagram names the \emph{real} agents shipped in
    \cmd{.claude/agents/} and \cmd{src/eda\_agents/templates/claude\_agents/}
    as of the current main: \cmd{analog-topology-recommender},
    \cmd{analog-sizing-advisor}, \cmd{digital-testbench-author},
    \cmd{gf180-docker-analog}, \cmd{gf180-docker-digital}. There is also a
    \cmd{librelane-expert} agent that specialises in config.yaml
    authoring.\par
    The dashed box above them is the MCP surface -- \cmd{eda-mcp} exposes
    skill templates, LUT access, and tool runners so any outer Claude
    Code session, or an ADK multi-agent hierarchy, can reuse the same
    primitives. We cover that in Section 5.\par
    Naming convention: \cmd{*-docker-*} agents drive
    \cmd{hpretl/iic-osic-tools} via \cmd{docker exec} and own the
    iteration inside the container. The others are planning / coding
    aids that run on the host.}
\end{frame}
